\scanpage{12}
$\operatorname{Spec} B$ sur $\operatorname{Spec} A$, et celles de $\operatorname{Spec} B_0$ sur $\operatorname{Spec} A_0$. \textcircled{b} Cette correspondance est \unsure{compatibles|compatible} avec les relations d'ordre, \deleted{\unsure{et en particulier}} \inserted{en particulier la relation} d'équivalence la plus (la moins) fine de \deleted{plus fine. On en conclut :} $\operatorname{Spec} B$ sur $\operatorname{Spec} A$, correspond à la relation d'équiv. la plus (la moins) fine de $\operatorname{Spec} B_0$ sur $\operatorname{Spec} A_0$.

[\,\textcircled{c} Si $J,J^{\unsure{0|o}}$ sont des idéaux d'\unsure{équivalences|équivalence} qui se correspondent, et $\deleted{\unsure{A|\Delta}} C$ ($C^0$) les ss-anneaux de $B$ ($B_0$) des invariants, alors $C^0\simeq C/IC\simeq C/I$, [de façon plus précise, si $\delta:B\longrightarrow D=B\otimes_A B/J$ et $\delta_0:B_0\longrightarrow D_0$ sont les homs \unsure{cobords|obvious}, alors on a
% slanted
\margin{expliciter plus \unsure{-\,-\,-\,-} \deleted{\unsure{équivalence ici} \unsure{de} $\unsure{ID|\delta|D}$}}
$\delta(b)\in ID$ sss $\deleted{b\in IB}\delta(b)=0$, i.e. $\delta_0(b_0)=0$ sss $\delta(b)=0$ i.e $b_0\in C_0$ sss $b\in C$], donc on écrira aussi $C_0$ aussi. Pour que $J$ soit effectif, il f et s. que $J^0$ le soit \ldots\ ]
% archive sheet number: 11
